\section{Post-Hoc Portfolio and Threshold Sensitivity}
\label{sec:equal_budget}

\subsection{Design and Scope}

The frozen graph action selects among six architectures with 24 trials in total, whereas the feature-only action contains one MLP with four. We examine the role of the candidate portfolio in a post-hoc reanalysis of the same records. Restricting the graph action to one architecture matches the number of trials at four per action. It does not equalize training time, FLOPs, preprocessing cost, or parameter count. The restriction changes both architecture availability and search breadth, so it does not isolate a causal effect of compute budget.

No model is retrained for the portfolio and threshold analyses below. In the portfolio restrictions, homophily and MLP validation accuracy remain unchanged, so the diagnostic takes the same actions as before. For each restricted portfolio we recompute the selected graph result, the margin-defined target, regret, and validation selection using the frozen selection and fallback rules. A wider candidate set guarantees a nondecreasing maximum validation accuracy; it does not guarantee a better selected test result. In the observed table, H2GCN alone has lower always-graph regret than the unrestricted portfolio, illustrating this distinction. The later preprocessing sensitivity in Section~\ref{sec:preprocessing_sensitivity} uses separately trained records.

The reconstruction script verifies all prospective record checksums against the release manifest, reuses the frozen assembly checks for configuration, provenance, trial selection, unique record scope, and splits, and regenerates the complete original audit before computing restrictions. Matching a few rounded aggregate values would not validate records for architectures that the full portfolio did not select. The unrestricted reconstruction recovers 89 graph targets and regrets of 0.26, 7.46, and 0.22 points for always-graph, Combined, and validation selection, respectively.

\subsection{Results}


\begin{table}[t]
\centering
\caption{Post-hoc single-architecture sensitivity at matched trial count. Each single-architecture graph action receives four trials against the MLP's four. Regret is mean full-set oracle-portfolio regret in percentage points. W/L/T counts datasets on which the combined rule attains lower, higher, or equal mean regret than always-graph.}
\label{tab:equal_budget}
\small
\begin{tabular}{@{}lrrrrrc@{}}
\toprule
Graph action & Trials & Graph targets & Always-graph & Combined & Validation & W/L/T \\
\midrule
Six architectures & 24 & 89/110 & 0.26 & 7.46 & 0.22 & 0/7/4 \\
\midrule
GAT & 4 & 41/110 & 6.35 & 2.08 & 0.01 & 4/2/5 \\
GCN & 4 & 46/110 & 5.29 & 1.74 & 0.01 & 3/3/5 \\
LINKX & 4 & 56/110 & 2.13 & 7.99 & 0.32 & 3/4/4 \\
GPR-GNN & 4 & 60/110 & 0.89 & 1.50 & 0.31 & 3/2/6 \\
GraphSAGE & 4 & 72/110 & 0.38 & 2.85 & 0.12 & 3/4/4 \\
H2GCN & 4 & 78/110 & 0.13 & 2.36 & 0.23 & 1/6/4 \\
\bottomrule
\end{tabular}
\end{table}


Table~\ref{tab:equal_budget} shows that the combined rule has lower mean full-set regret than always-graph under GCN and GAT, reversing its relative performance in the full-portfolio comparison. The difference is substantial descriptively: 1.74 versus 5.29 points for GCN and 2.08 versus 6.35 for GAT. Under LINKX the ordering reverses, at 7.99 versus 2.13 points. GPR-GNN, GraphSAGE, and H2GCN also favor always-graph on mean regret. Thus the same diagnostic has different measured value against the same trivial policy when the graph action changes.

The complete table does not support a simple ordering of architecture categories. GPR-GNN, one of the heterophily-aware architectures, yields the lowest absolute Combined regret (1.50 points), below GCN (1.74) and GAT (2.08). Absolute regret and improvement over always-graph answer different questions. Each row also uses its own graph/MLP oracle, so absolute regrets across rows are not losses relative to one common oracle. The evidence supports variation with the candidate portfolio, rather than a claim that a broad architecture class necessarily makes diagnostics effective or ineffective.

Dataset-level heterogeneity remains substantial. For GCN, Combined wins, loses, and ties always-graph on three, three, and five datasets; for GAT the counts are four, two, and five. Those counts and the reported means are descriptive. The original eight-comparison primary-analysis family did not include these six restrictions, and these data were already available when the additional analysis was formulated. We do not infer equivalence, a stable out-of-dataset advantage, or a lack of diagnostic information from the small number of informative datasets.

\subsection{Interpretation}

The restriction analysis provides direct empirical evidence that evaluation against one graph action need not transfer to another. It is consistent with a mismatch between a fixed homophily rule and some candidate portfolios, but it does not prove a universal mechanism or exclude other threshold choices. The operational lesson is to specify the candidate models and resource budget when evaluating a diagnostic, and to test any proposed improvement on independent data.

\subsection{All Candidate Subsets and Dataset Sensitivity}
\label{sec:subset_robustness}

We enumerate all $2^6-1=63$ nonempty graph-architecture subsets using the same validation selection and frozen diagnostic actions. Define $\Delta$ as Combined minus always-graph mean regret in percentage points; negative values favor Combined. Table~\ref{tab:subset_sensitivity} reports the entire enumeration by subset size. Only GAT, GCN, and their pair favor Combined. The 63 subsets overlap in both models and datasets: their counts are descriptive, not 63 independent replications or estimates of a generalization probability.

\begin{table}[t]
\centering
\caption{Post-hoc enumeration of all available graph portfolios. Each architecture contributes four trials. The range is over candidate subsets, not a confidence interval.}
\label{tab:subset_sensitivity}
\small
\begin{tabular}{@{}rrrrr@{}}
\toprule
Architectures & Subsets & $\Delta<0$ & Min $\Delta$ & Max $\Delta$ \\
\midrule
1 & 6 & 2 & -4.26 & 5.86 \\
2 & 15 & 1 & -3.23 & 7.32 \\
3 & 20 & 0 & 1.68 & 7.38 \\
4 & 15 & 0 & 2.48 & 7.38 \\
5 & 6 & 0 & 3.09 & 7.31 \\
6 & 1 & 0 & 7.20 & 7.20 \\
\bottomrule
\end{tabular}
\end{table}

At the same eight-trial graph budget, GAT+GCN yields Combined/always-graph regrets of 2.27/5.50 points, whereas H2GCN+LINKX yields 7.59/0.27. Thus equal trial count does not remove variation with architecture composition. Nevertheless, the favorable average under GAT or GCN is sensitive to dataset composition. Omitting any one dataset preserves its direction, but omitting both Cornell and Wisconsin reverses it: $\Delta$ becomes $+1.11$ for GAT and $+0.92$ for GCN, and none of the 63 subsets favors Combined. This is a post-hoc influence check, not grounds for discarding these datasets.

Separately omitting Chameleon and Squirrel, whose original versions have known duplicate-node concerns, leaves full-portfolio Combined/always-graph regrets of 4.90/0.31 points ($\Delta=+4.59$). Seven of the 63 remaining subset comparisons favor Combined. This exclusion does not establish performance on cleaned versions of either graph. For the full portfolio, omitting any single dataset leaves Combined-minus-always-graph regret between 5.556 and 7.922 points, without reversing its sign. All per-subset and per-dataset values, including leave-one-dataset-out ranges, are released; neither exclusion replaces the frozen full-set result.

\paragraph{Fallback and dataset-exclusion sensitivity.}
Table~\ref{tab:fallback_sensitivity} crosses MLP versus graph fallback with inclusion versus exclusion of Chameleon and Squirrel, holding the six-architecture portfolio and all recorded model outcomes fixed. Combined has higher descriptive mean regret than always-graph in each of these four configurations, but the excess varies from 0.141 to 7.201 points. This range applies only to these four configurations: the previously reported Cornell/Wisconsin exclusion gives an 8.683-point excess under MLP fallback. Fallback and dataset groupings overlap, so they do not supply independent causal contributions. Removing datasets with known duplicate-node concerns does not identify those duplicates as the cause of the losses. Using the frozen validation-selection rule only for abstentions instead gives 1.806 points on all datasets and 0.442 after excluding Chameleon/Squirrel, close to graph fallback; explicit graph and MLP decisions are unchanged.

\begin{table}[ht]
\centering
\caption{Post-hoc fallback and dataset-exclusion sensitivity of the frozen six-architecture portfolio. Regrets and Combined-minus-always-graph differences are percentage points. All entries reuse recorded outcomes; no model is retrained.}
\label{tab:fallback_sensitivity}
\small
\begin{tabular}{@{}llrrr@{}}
\toprule
Datasets & Fallback & Combined & Always-graph & Difference \\
\midrule
All 11 & MLP & 7.457 & 0.256 & 7.201 \\
All 11 & Graph & 1.815 & 0.256 & 1.559 \\
Exclude Chameleon/Squirrel & MLP & 4.898 & 0.313 & 4.586 \\
Exclude Chameleon/Squirrel & Graph & 0.454 & 0.313 & 0.141 \\
\bottomrule
\end{tabular}
\end{table}

\paragraph{Headroom and an optimistic threshold bound.}
Always-graph's regret is also the maximum possible reduction in mean regret relative to that reference, holding each portfolio's recorded outcomes and oracle fixed. This headroom is only 0.256 points for the full portfolio, versus 5.292 for GCN and 6.346 for GAT. Exhaustively evaluating all distinct decisions of graph-if-$h_1\geq t$ for $t\in[0,1]$ gives a test-outcome-selected minimum of 0.256, 0.812, and 0.678 points, respectively. Thus no cutoff in this one-sided family improves on always-graph in the full-portfolio records; this is not a result about every diagnostic or an external dataset. For GCN and GAT, the frozen rule closes 67.04\% and 67.17\% of the respective headroom, while the optimistic cutoff bound closes 84.67\% and 89.32\%; validation selection closes about 99.9\%. Appendix~\ref{app:threshold_headroom} reports every single architecture and the exact curves. These bounds reuse test outcomes to choose the cutoff and are not calibrated or deployable performance estimates. The held-out analysis below retains its original grid and results.

\subsection{A Post-Hoc Calibrated Threshold Control}
\label{sec:calibrated_threshold}

To distinguish the frozen threshold from the homophily statistic itself, we fit a graph-if-$h_1\geq t$ rule separately for each held-out dataset. The fixed grid is $\{-1,0,0.01,\ldots,1,2\}$, including constant graph and MLP policies. Each fold minimizes equal-weight mean dataset regret on the other ten datasets; ties within $10^{-12}$ choose the smallest threshold. The selected-model test outcomes of those ten datasets serve as meta-training targets. Neither outcomes nor model records of the held-out dataset select its threshold. We then evaluate on that dataset's ten seeds and average the eleven held-out dataset means.

\begin{table}[t]
\centering
\caption{Post-hoc held-out threshold control. Regrets are percentage points, using each row's own graph/MLP oracle. Calibration is performed separately for each portfolio; it is not an independent prospective experiment.}
\label{tab:calibrated_threshold}
\small
\begin{tabular}{@{}lrrr@{}}
\toprule
Graph action & Frozen Combined & Calibrated & Always-graph \\
\midrule
GAT & 2.08 & 1.53 & 6.35 \\
GCN & 1.74 & 1.58 & 5.29 \\
GPR-GNN & 1.50 & 0.81 & 0.89 \\
GraphSAGE & 2.85 & 0.82 & 0.38 \\
H2GCN & 2.36 & 1.05 & 0.13 \\
LINKX & 7.99 & 4.15 & 2.13 \\
Six architectures & 7.46 & 2.41 & 0.26 \\
\bottomrule
\end{tabular}
\end{table}

Calibration lowers mean regret relative to the frozen rule in all seven rows of Table~\ref{tab:calibrated_threshold}, including from 7.46 to 2.41 points for the full portfolio, but still exceeds its always-graph regret of 0.26. In that full-portfolio control, 10 of the 11 folds select the constant always-graph policy ($t=-1$); the Roman-empire fold selects $t=0.11$ and contributes approximately 89\% of the remaining aggregate regret. The remaining loss is therefore concentrated in one held-out dataset. Calibration retains lower descriptive regret than always-graph for GCN and GAT and also improves slightly on it for GPR-GNN, unlike the frozen threshold. These results argue against interpreting failure of the frozen cutoff as failure of every threshold. The control was designed after inspection of the original benchmark, its training folds overlap, and graph domains may be related. It is a stronger threshold baseline, not a representative comparison with all published diagnostic methods or evidence of stable external validity. These post-hoc comparisons are descriptive and do not enter the pre-specified primary test family.

\subsection{Feature Preprocessing Sensitivity}
\label{sec:preprocessing_sensitivity}

The frozen run applies PyG 2.7.0 \texttt{NormalizeFeatures} to every dataset on load. This transform subtracts the global minimum of the feature matrix and divides each row by its sum clamped below at one. This common choice is not a neutral operation for the continuous features in the Heterophilous Graph Dataset. We therefore retrained all seven architectures for ten seeds on Roman-empire and Amazon-ratings under the frozen transform and under raw features, keeping the graph, splits, seeds, trial grid, checkpoint rule, and device paired. This is a post-hoc two-dataset sensitivity experiment; it does not replace the frozen 11-dataset benchmark or support a population-wide preprocessing claim.

\begin{table}[t]
\centering
\caption{Post-hoc paired feature-preprocessing sensitivity on two heterophilous datasets. Values are mean test accuracy or full-set oracle regret in percentage points; graph targets are out of ten seeds. Each graph action still searches six architectures and 24 trials.}
\label{tab:preprocessing_sensitivity}
\small
\begin{tabular}{@{}llrrrrr@{}}
\toprule
Dataset & Features & MLP & Graph & Graph--MLP & Combined regret & Graph targets \\
\midrule
Roman-empire & NormalizeFeatures & 16.71 & 42.16 & 25.45 & 25.45 & 10/10 \\
Roman-empire & Raw & 66.19 & 80.96 & 14.77 & 14.77 & 10/10 \\
Amazon-ratings & NormalizeFeatures & 36.76 & 53.06 & 16.30 & 16.30 & 10/10 \\
Amazon-ratings & Raw & 41.35 & 53.17 & 11.82 & 11.82 & 10/10 \\
\bottomrule
\end{tabular}
\end{table}

Raw features raise selected MLP accuracy by 49.48 points on Roman-empire and 4.59 points on Amazon-ratings, while the selected graph changes by 38.81 and 0.11 points, respectively. The graph portfolio remains the better action in all 20 paired units under both conditions, so always-graph and validation selection have zero regret in this restricted sensitivity set; the homophily rule still selects MLP because both datasets lie below $h_1=0.55$. The result shows that the frozen preprocessing choice materially changes the graph--MLP estimand on these two datasets. It does not show that normalization caused the full benchmark's negative result, because the comparison covers only two selected datasets and remains post-hoc.

The low normalized-feature MLP results also warrant checking training behavior. Appendix~\ref{app:training_diagnostics} reports separate validation-only input-parameterization experiments and a bounded training-repeatability audit. Centering and scalar scaling of normalized features recover MLP validation performance in the tested conditions, while repeated LINKX training under identical inputs can diverge substantially. The separately retrained normalized-feature condition is a sensitivity reference, not an exact replay of the original selected models: its graph--MLP gaps are 25.45 versus 23.66 points in the frozen run for Roman-empire, and 16.30 versus 16.45 for Amazon-ratings. Accordingly, Table~\ref{tab:preprocessing_sensitivity} preserves the recorded paired test results; it is not a guarantee of repeatable graph effect sizes across executions. Neither these two datasets nor the later validation-only diagnostics establish that the complete 11-dataset result is robust to alternative preprocessing.

\paragraph{Fallback within the paired preprocessing run.}
Table~\ref{tab:preprocessing_fallback} recomputes the two fallback choices from the retained per-seed results of this two-dataset experiment. All 20 normalized-feature units abstain, whereas all 20 raw-feature units explicitly select MLP because each selected MLP validation accuracy exceeds 0.40. Consequently, graph fallback changes only the normalized-feature results. This interaction describes alternative evaluation configurations; a validation score below 0.40 is not by itself evidence that graph will outperform MLP. The table neither pools these later model outcomes with the original 11-dataset records nor defines a corrected benchmark effect.

\begin{table}[ht]
\centering
\caption{Post-hoc fallback sensitivity within the same paired preprocessing experiment on Roman-empire and Amazon-ratings (20 units per feature condition). Entries are mean full-set regret in percentage points. The graph action uses all six architectures.}
\label{tab:preprocessing_fallback}
\small
\begin{tabular}{@{}lrrr@{}}
\toprule
Features & Combined, MLP fallback & Combined, graph fallback & Always-graph \\
\midrule
NormalizeFeatures & 20.876 & 0.000 & 0.000 \\
Raw & 13.299 & 13.299 & 0.000 \\
\bottomrule
\end{tabular}
\end{table}
